% Banco de dados
with app.app_context():
    database.drop_all()
    database.create_all()
    senha_hash = generate_password_hash('kauan123')
    usuario = Usuario(nome_usuario='kauan medeiros',email='kauanmm11.web@gmail.com',senha=senha_hash,tipo='admin')
    database.session.add(usuario)
    database.session.commit()
% 

% css

<link rel="stylesheet" href="{{ url_for('static', filename='css/base.css') }}">

%

% Flask inicio dentro __init__py
from flask import Flask
from flask_sqlalchemy import SQLAlchemy

app = Flask(__name__)
app.config['SECRET_KEY'] = '416e7e3da3790330038b0a7f997ef687'
app.config['SQLALCHEMY_DATABASE_URI'] = "sqlite:///app.db"

database = SQLAlchemy(app)

from kauanmedeirosportfolio.auth import home, logou
from kauanmedeirosportfolio import routes
% 

% form.py -> import
from flask_wtf import FlaskForm
from flask_wtf.file import FileField, FileAllowed
from wtforms import StringField, PasswordField, EmailField, SubmitField ,BooleanField, URLField
from wtforms.validators import Email, EqualTo, Length, DataRequired

criar categoria
categoria = SelectField(
'Categoria',choices=[('Desenvolvimento Web', 'Desenvolvimento Web'),('Dashboard', 'Dashboard'),('Landing Page','Landing Page'),('Site Institucional', 'Site Institucional')],validators=[DataRequired()])
% 

% admin -> decorators.py
def admin_required(func):
    @wraps(func)
    def wrapper(*args, **kwargs):

        if 'usuario_id' not in session:
           flash('Usuario não esta logado, faça o login e tente novamente')
           return redirect(url_for('home'))

        usuario = Usuario.query.filter_by(id=session['usuario_id']
        ).first()

        if usuario is None:
            flash('Usuario não esta logado, faça o login e tente novamente')
            return redirect(url_for('home'))

        if usuario.tipo != 'admin':
            flash('Usuario não permidido')
            return redirect(url_for('home'))

        return func(*args, **kwargs)
    return wrapper
% 

% models.py banco de dados
class Usuario(database.Model):
    id = database.Column(database.Integer(), primary_key=True)
    nome_usuario = database.Column(database.String(), nullable=False)
    email = database.Column(database.String(), nullable=False, unique=True)
    senha = database.Column(database.String(), nullable=False)
    tipo = database.Column(database.String(), default='cliente')

    entre ligar usuario com outra coisa no exemplo esta no posts
    usuario = database.relationship(Usuario,backref='posts')
    id_usuario = database.Column(database.Integer(),database.ForeignKey('usuario.id'))

    crir categoria
    categoria = database.Column(database.String(100), nullable=False)
% 

% pasta -> templates e static
nome do projeto e colocar tudo dentro menos main.py
routes, css, js e img
% 